\chapter{المجاميع الأسية وطريقة فان دير كوربوت}
\label{ch:exponential-sums-van-der-corput}

\section*{نطاق الفصل وحالته}

يبني هذا الفصل الآلة الكلاسيكية التي تحوّل تذبذب الطور إلى إلغاء كمي في المجاميع الأسية. يثبت داخليًا الحد التافه، وصيغة الجمع الجزئي المنفصلة، ومتباينة فرق فان دير كوربوت، واختبار المشتقة الأولى في الصيغة المضبوطة التي تقيس البعد عن الأعداد الصحيحة. أما اختبار المشتقة الثانية، وعملية \(B\)، والإطار العام للأزواج الأسية، فتعرض بوصفها نتائج مقتبسة ومشروحة. جميع نتائج هذا الفصل في هذه المرحلة \textbf{مسودة غير قابلة للاستشهاد} إلى أن يكتمل البناء والتدقيق الرياضي والمرجعي والمراجعة المستقلة.

\subsection*{نواتج التعلم}

بنهاية الفصل يستطيع القارئ أن:
\begin{enumerate}
\item يفسر الفرق بين الحد التافه والإلغاء الحقيقي في مجموع أسي.
\item ينقل تقديرًا لمجموعات جزئية غير موزونة إلى مجموع موزون باستعمال الجمع الجزئي.
\item يشتق متباينة فرق فان دير كوربوت ويفهم دور الارتباطات المزاحة.
\item يميز بين الشرط الصحيح \(\|f'(x)\|\ge\lambda\) والشرط غير الكافي \(|f'(x)|\ge\lambda\).
\item يشرح الفرق البنيوي بين عمليتي \(A\) و\(B\)، وحدود ما يثبته الفصل داخليًا.
\item يربط الأدوات العامة هنا بحاجات الأقواس الصغرى في الفصل السابق من دون اعتماد دائري.
\end{enumerate}

\section{المجموع الأسي ومعنى الإلغاء}

نستخدم طوال الفصل التطبيع
\[
e(t)=e^{2\pi i t}.
\]
إذا كانت \(I\subset\mathbb R\) فترة محدودة و\(f:I\to\mathbb R\)، نكتب
\[
S(f;I)=\sum_{n\in I\cap\mathbb Z} e(f(n)).
\]
وفي النسخة الموزونة نكتب
\[
S(a,f;I)=\sum_{n\in I\cap\mathbb Z} a_n e(f(n)).
\]
كل حد في المجموع الموزون له طول \(|a_n|\)، ولذلك ينتج فورًا الحد التالي.

\begin{proposition}[الحد التافه]
\label{prop:exp-trivial-bound}
لكل متتالية مركبة منتهية \((a_n)\) ولكل طور حقيقي \(f\)،
\[
\left|\sum_n a_n e(f(n))\right|
\le
\sum_n |a_n|.
\]
وبخاصة، إذا كان عدد الحدود \(N\) وكانت \(|a_n|\le1\)، فإن حجم المجموع لا يتجاوز \(N\).
\resultid{ANT-ID-18-01}
\provedhere
\end{proposition}

\begin{proof}
هذه هي متباينة المثلث مباشرة، لأن \(|e(f(n))|=1\):
\[
\left|\sum_n a_n e(f(n))\right|
\le
\sum_n |a_n e(f(n))|
=
\sum_n|a_n|.
\]
\end{proof}

الغاية من نظرية المجاميع الأسية ليست إعادة إنتاج هذا الحد، بل إثبات أن اتجاهات الحدود المركبة تتغير بما يكفي كي يلغي بعضها بعضًا. عندما نحصل مثلًا على
\[
S(f;I)\ll |I|^{1-\delta}
\qquad(\delta>0),
\]
نقول إن هناك \emph{ادخارًا بقوة} عن الحد التافه.

\begin{example}[طور خطي]
إذا كان \(f(n)=\alpha n+\beta\)، فإن
\[
\sum_{M<n\le M+N}e(f(n))
=e(\beta)e(\alpha(M+1))
\frac{1-e(\alpha N)}{1-e(\alpha)}
\]
عندما \(\alpha\notin\mathbb Z\). ومن ثم
\[
\left|\sum_{M<n\le M+N}e(\alpha n+\beta)\right|
\le
\min\!\left(N,\frac{1}{2\|\alpha\|}\right),
\]
حيث \(\|x\|\) هي المسافة من \(x\) إلى أقرب عدد صحيح. تكشف هذه الصيغة منذ البداية أن الكمية الصحيحة ليست \(|\alpha|\)، بل \(\|\alpha\|\).
\end{example}

\section{الجمع الجزئي ونقل الأوزان}

كثيرًا ما نملك تقديرًا للمجاميع الجزئية
\[
A(x)=\sum_{m<n\le x}a_n,
\]
ثم نحتاج إلى إدخال وزن يتغير ببطء. الأداة المناسبة هي النسخة المنفصلة من التكامل بالتجزئة.

\begin{lemma}[الجمع الجزئي المنفصل]
\label{lem:exp-partial-summation}
لتكن \(a<n\le b\) أعدادًا صحيحة، ولتكن \(u_n\) متتالية مركبة و\(g\in C^1([a,b])\). ضع
\[
U(x)=\sum_{a<n\le x}u_n.
\]
عندئذ
\[
\sum_{a<n\le b}u_n g(n)
=U(b)g(b)-\int_a^b U(t)g'(t)\,dt.
\]
وبالتالي
\[
\left|\sum_{a<n\le b}u_n g(n)\right|
\le
\sup_{a\le t\le b}|U(t)|
\left(|g(b)|+\int_a^b|g'(t)|\,dt\right).
\]
\resultid{ANT-LEM-18-01}
\provedhere
\end{lemma}

\begin{proof}
نكتب \(u_n=U(n)-U(n-1)\)، ثم نجمع بالتلسكوب:
\[
\sum_{a<n\le b}u_ng(n)
=U(b)g(b)-\sum_{a<n<b}U(n)\bigl(g(n+1)-g(n)\bigr).
\]
وباستخدام
\[
g(n+1)-g(n)=\int_n^{n+1}g'(t)\,dt
\]
وتعريف \(U(t)=U(n)\) على \([n,n+1)\)، نحصل على الصيغة التكاملية. أما المتباينة فتنتج بأخذ القيم المطلقة.
\end{proof}

\begin{remark}
هذه اللمّة لا تولّد الإلغاء بذاتها؛ إنها تحفظ الإلغاء الموجود في المجاميع الجزئية وتنقله إلى أوزان ملساء أو رتيبة مع خسارة تقاس بتغير الوزن الكلي.
\end{remark}

\section{متباينة فرق فان دير كوربوت}

الفكرة الجوهرية هي أن المجموع الأصلي قد يكون صعبًا، بينما تكون ارتباطاته مع نسخ مزاحة منه أسهل. لتكن
\[
Z=\sum_{n=1}^{N}z_n,
\]
ونمدد \(z_n\) بالصفر خارج \(1\le n\le N\). لكل \(1\le H\le N\)، لدينا
\[
H Z
=
\sum_{m=1}^{N+H-1}\sum_{h=0}^{H-1}z_{m-h}.
\]
تطبيق كوشي--شفارتس على المجموع الخارجي يقود إلى المتباينة الأساسية.

\begin{lemma}[متباينة فرق فان دير كوربوت]
\label{lem:van-der-corput-differencing}
إذا كان \(z_1,\dots,z_N\in\mathbb C\) و\(1\le H\le N\)، فإن
\[
\left|\sum_{n=1}^{N}z_n\right|^2
\le
\frac{N+H-1}{H}
\left(
\sum_{n=1}^{N}|z_n|^2
+2\sum_{h=1}^{H-1}
\left(1-\frac{h}{H}\right)
\left|\sum_{n=1}^{N-h}z_{n+h}\overline{z_n}\right|
\right).
\]
\resultid{ANT-LEM-18-02}
\provedhere
\end{lemma}

\begin{proof}
بتمديد \(z_n\) بالصفر خارج المجال نحصل على
\[
H\sum_{n=1}^{N}z_n
=
\sum_{m=1}^{N+H-1}\sum_{h=0}^{H-1}z_{m-h}.
\]
ومن كوشي--شفارتس،
\[
H^2|Z|^2
\le
(N+H-1)
\sum_m\left|\sum_{h=0}^{H-1}z_{m-h}\right|^2.
\]
نوسع المربع ونرتب الحدود بحسب الفرق \(r=h-k\). الحد القطري \(r=0\) يظهر \(H\) مرة، وكل فرق موجب \(1\le r\le H-1\) يظهر \(H-r\) مرة مع مرافقه المركب. لذلك
\[
\sum_m\left|\sum_{h=0}^{H-1}z_{m-h}\right|^2
=
H\sum_{n=1}^{N}|z_n|^2
+2\operatorname{Re}\sum_{r=1}^{H-1}(H-r)
\sum_{n=1}^{N-r}z_{n+r}\overline{z_n}.
\]
نستبدل الجزء الحقيقي بالقيمة المطلقة، ثم نقسم على \(H^2\)، فنحصل على الصيغة المطلوبة.
\end{proof}

\begin{corollary}[نسخة الأطوار]
\label{cor:van-der-corput-phase}
إذا كان \(z_n=e(f(n))\)، فإن
\[
|S(f;[1,N])|^2
\le
\frac{N+H-1}{H}
\left(
N+2\sum_{h=1}^{H-1}
\left(1-\frac{h}{H}\right)
\left|\sum_{n=1}^{N-h}e\bigl(f(n+h)-f(n)\bigr)\right|
\right).
\]
\end{corollary}

\begin{proof}
لدينا \(|z_n|=1\)، كما أن
\[
z_{n+h}\overline{z_n}
=e\bigl(f(n+h)-f(n)\bigr).
\]
نعوض في اللمّة السابقة.
\end{proof}

\begin{remark}[اختبار الثابت]
إذا أخذنا \(z_n=1\)، فإن طرف الارتباط يساوي \(N-h\)، ولا تنتج المتباينة ادخارًا وهميًا. وهذا متوقع لأن الطور الثابت لا يملك أي تذبذب.
\end{remark}

\section{من الفروق إلى المشتقات}

عند تطبيق المتباينة على طور أملس، يصبح الطور الجديد
\[
\Delta_hf(x)=f(x+h)-f(x).
\]
ومن مبرهنة القيمة المتوسطة،
\[
\Delta_hf(x)=h f'(\xi_x)
\]
لنقطة مناسبة \(\xi_x\in(x,x+h)\). وهكذا تخفض عملية الفرق رتبة التعقيد المشتقي: المعلومات عن \(f''\) تتحول إلى معلومات عن مشتقة \(\Delta_hf\)، والمعلومات عن مشتقات أعلى تنتقل بالطريقة نفسها. هذه هي النواة التحليلية لعملية \(A\) في طريقة الأزواج الأسية.

لكن يجب الحذر: إذا كان \(f'(x)\) قريبًا من عدد صحيح، فقد تكون الحدود \(e(f(n))\) شبه متراصفة حتى لو كانت \(|f'(x)|\) كبيرة. لذلك سيكون اختبار المشتقة الأولى في القسم التالي مبنيًا على الشرط
\[
\|f'(x)\|\ge\lambda>0,
\]
لا على \(|f'(x)|\ge\lambda\) وحده.

\section*{حالة دفعة التأليف الأولى}

\begin{center}
\begin{tabular}{ll}
\texttt{ANT-ID-18-01} & \texttt{AUTHORED-DRAFT / PROVED-HERE}\\
\texttt{ANT-LEM-18-01} & \texttt{AUTHORED-DRAFT / PROVED-HERE}\\
\texttt{ANT-LEM-18-02} & \texttt{AUTHORED-DRAFT / PROVED-HERE}\\
\texttt{ANT-THM-18-01} & \texttt{RESERVED / NEXT BATCH}\\
\texttt{ANT-THM-18-02} & \texttt{RESERVED / CITED-EXPLAINED}\\
\texttt{ANT-DEF-18-01} & \texttt{RESERVED / CITED-FRAMEWORK}\\
\texttt{ANT-PROP-18-01} & \texttt{RESERVED / LIMITED INTERNAL VERSION}\\
\texttt{ANT-PROP-18-02} & \texttt{RESERVED / CITED-EXPLAINED}
\end{tabular}
\end{center}

لا تمثل هذه الحالات اعتمادًا نهائيًا؛ يبقى الفصل \texttt{AUTHORED-DRAFT / NON-CITABLE} حتى اكتمال بقية الوحدات والبناء والمراجعة المستقلة.
